\begin{figure}[htbp]
\centering
\begin{tikzpicture}
\node[booknode,text width=3cm] (a) at (1.6,2.6) {1. A0 取得 S1};
\node[booknode,text width=3cm] (b) at (5.6,2.6) {2. A0 准备成功};
\node[booknode,text width=3cm] (c) at (9.6,2.6) {3. A1 也取得 S1};
\node[bookresult,text width=3cm] (d) at (1.6,.4) {4. A0 提交\\余额 7000};
\node[booknode,text width=3cm] (e) at (5.6,.4) {5. A1 准备成功};
\node[bookresult,text width=3cm] (f) at (9.6,.4) {6. A1 提交\\余额 4000};
\draw[bookarrow] (a) -- (b);
\draw[bookarrow] (b) -- (c);
\draw[bookarrow] (c.south) -- (9.6,1.5) -- (1.6,1.5) -- (d.north);
\draw[bookarrow] (d) -- (e);
\draw[bookalert] (e) -- (f);
\node[booknote,text width=10.5cm] at (5.6,-1) {账户排队与余额条件仍成立；删除完整 S 排他后，同一身份形成两份成功};
\end{tikzpicture}
\caption{去掉身份唯一性后的六步反例}
\label{fig:ch06-no-unique}
\end{figure}
